\scanpage{201}%
\mspage{75}%
\deleted{\unsure{tion}} d) (p. 45). De $=$ la condition d'exactitude e) (p. 46) sur les $\underrightarrow{\lim}$ filtrantes (\unsure{dérivant|devient} \unsure{d'un|donc} \unsure{lemme|jeu} \unsure{connu} \unsure{en moins|au moins}).

Pour être tranquille, je vais faire sur $M$ les hypothèses a) à e) inclus. (bien que e) me semble un peu artificielle, on devrait pouvoir l'éliminer, dans un second souffle !) Ces conditions ont donc les propriétés de stabilité souhaitables par passage aux $M(S)$, et en particulier par passage à $\mathrm{Fl}(M)$. Quitte à augmenter $\unsure{\pi|R}_0$, on peut donc supposer que les
\[
 \mathrm{Fl}^{\pi}M=\mathrm{Hom}\,(\Delta^1,\deleted{\unsure{\mathrm{Fl}}}M^{\pi})
 \qquad(\pi\geq\pi_0)
\]
forment aussi une filtration cardinale de $\underline{\mathrm{Fl}}\,M$. \deleted{\illegible{1}}

On dispose donc des notions de parties $\pi$-accessibles, accessibles, L-$\pi$-accessibles, L-accessibles de $\mathrm{Fl}\,M=\mathrm{Ob}\,\underline{\mathrm{Fl}}\,(M)$. \deleted{\unsure{Cela|Ces}} \unsure{les|Ces} notions sont stables par réunions et intersections finies, et les notions absolues (sans $\pi$) stables par intersections ``quelconques'' (i.e. \deleted{\unsure{pour une}} \inserted{pour} des \deleted{\unsure{cardin}} \uline{petites} familles de parties). Ce sera là \unsure{un des cas|une de mes} parti%
% archive sheet number: 200
